\chapter{Método Generalizado dos Momentos}
\label{cap:gmm}
% ── índice ──────────────────────────────────────────────────────
\index{método dos momentos generalizados|textbf}
\index{GMM|see{método dos momentos generalizados}}
\index{gmm@\texttt{gmm()}|textbf}
\index{sobreidentificação|textbf}
\index{teste J|textbf}
\index{condições de momento}

% ─────────────────────────────────────────────────────────────────────────────
\section{O framework}

O Método Generalizado dos Momentos (GMM) unifica os estimadores dos capítulos
anteriores sob um único princípio: a identificação via \textbf{condições de
momento} populacionais.

Um modelo GMM especifica $L$ condições de momento:
\[
  E\!\left[g_i(\beta)\right] = 0, \qquad g_i(\beta) \in \mathbb{R}^L
\]

OLS e IV são casos especiais:

\begin{center}
\begin{tabular}{lll}
\toprule
\textbf{Estimador} & \textbf{Condição de momento} & \textbf{Situação} \\
\midrule
OLS   & $E[x_i\,\varepsilon_i] = 0$ & $L = K$, identificado \\
IV/2SLS & $E[z_i\,\varepsilon_i] = 0$ & $L = K$, identificado \\
GMM   & $E[z_i\,\varepsilon_i] = 0$ & $L > K$, sobreidentificado \\
\bottomrule
\end{tabular}
\end{center}

Com $L = K$ momentos, resolve-se $g(\hat\beta) = 0$ exatamente — equivalente
ao IV just-identified. Com $L > K$, o sistema é sobredeterminado e não existe
$\beta$ que zere todos os momentos simultaneamente; o GMM minimiza uma forma
quadrática ponderada:

\[
  \hat\beta_{\mathrm{GMM}}
  = \arg\min_\beta\; n\, \bar g(\beta)' \hat W^{-1} \bar g(\beta)
\]

onde $\bar g(\beta) = \tfrac{1}{n}\sum_i g_i(\beta)$ e $\hat W$ é a estimativa
consistente da variância assintótica de $\sqrt{n}\,\bar g(\beta)$.

% ─────────────────────────────────────────────────────────────────────────────
\section{O estimador dois passos}

\begin{enumerate}
  \item \textbf{Primeiro passo:} estima $\hat\beta_1$ com $W = I$ (peso
        uniforme).
  \item \textbf{Matriz de peso eficiente:} calcula $\hat W = \tfrac{1}{n}
        \sum_i z_i z_i' \hat\varepsilon_i^2$ com resíduos do primeiro passo.
  \item \textbf{Segundo passo:} reestima $\hat\beta_2$ com $\hat W^{-1}$ —
        o estimador GMM eficiente.
\end{enumerate}

O estimador dois passos é assintoticamente eficiente na classe de estimadores
GMM com as mesmas condições de momento.

% ─────────────────────────────────────────────────────────────────────────────
\section{Teste J de sobreidentificação}

Quando $L > K$, é possível testar se os momentos excedentes são válidos.
A estatística de Hansen (\emph{J-test}):

\[
  J = n\, \bar g(\hat\beta)' \hat W^{-1} \bar g(\hat\beta)
  \;\xrightarrow{d}\; \chi^2(L - K)
  \quad \text{sob } H_0: E[g_i(\beta_0)] = 0
\]

$H_0$ é a hipótese de que \emph{todos} os $L$ instrumentos são exógenos.
Rejeitar $H_0$ indica que pelo menos um instrumento é correlacionado com
$\varepsilon$ — o modelo está mal especificado ou um instrumento é inválido.

\begin{aviso}
  O teste J é necessário mas não suficiente: não rejeitar não prova que todos
  os instrumentos são válidos, apenas que as restrições de sobreidentificação
  são consistentes com os dados.
\end{aviso}

% ─────────────────────────────────────────────────────────────────────────────
\section{Verificação por DGP}

\subsection*{OLS, IV e GMM lado a lado}

O DGP repete a estrutura do capítulo~\ref{cap:iv}: $x$ é endógeno ($\mathrm{Cov}(x,\varepsilon)
\neq 0$), com dois instrumentos válidos $z_1$ e $z_2$. O $\varepsilon$ é
ortogonalizado contra $z_1$ in-sample — IV just-identified recupera
exatamente; GMM com dois instrumentos é sobreidentificado ($L=2, K=1$).

\begin{lstlisting}[caption={DGP GMM — OLS viesado, IV exato, GMM eficiente}]
seed(42)
input df
id
1
end
for i in 1..=10 { df = append(df, df) }
df |> filter(_n <= 1000)
generate df z1  = rnormal()
generate df z2  = rnormal()
generate df v   = rnormal()
let d1          = cov(df, z1, v) / variance(df, z1)
let d0          = mean(df, v) - d1 * mean(df, z1)
generate df eps = v - d0 - d1*z1
generate df x   = 0.8*z1 + 0.4*z2 + eps
generate df y   = 1.5 + 2.0*x + 0.7*eps
let m_ols = ols(y ~ x, df)
let m_iv  = iv(y ~ x, ~ z1, df)
let m_gmm = gmm(y ~ x, ~ z1 + z2, df)
display m_ols
display m_iv
display m_gmm
\end{lstlisting}

\begin{lstlisting}[language={}, basicstyle=\ttfamily\scriptsize,
  frame=none, numbers=none, backgroundcolor=\color{codbg},
  xleftmargin=0pt, xrightmargin=0pt, breaklines=false]
Dep. Variable: y  R²=0.9798  N=1000  F=48418.53
const   1.2757***  (0.0077)    x   2.3736***  (0.0108)

Dep. Variable: y  Estimator: 2SLS  N=1000
const   1.5000***  (0.0170)    x   2.0000***  (0.0264)

Dep. Variable: y  GMM Two-Step  N=1000  J=2.5012  Prob(J)=0.1138  DF=1
x0   1.4875***  (0.0145)    x1   2.0206***  (0.0226)
\end{lstlisting}

OLS é viesado (endogeneidade): $\hat\beta_1 = 2{,}37 \neq 2{,}0$. IV com $z_1$
recupera exatamente $2{,}0000$ (ortogonalização in-sample). GMM com $z_1 + z_2$
converge a $2{,}0206$ — consistente, não exato, porque usa dois momentos para
estimar um parâmetro. O teste J ($J = 2{,}50$, $p = 0{,}11$) não rejeita
$H_0$: ambos os instrumentos são válidos.

\begin{nota}
  O Hayashi exibe os coeficientes como \texttt{x0} (constante) e \texttt{x1}
  (slope) na saída do GMM. Isso é uma particularidade do display interno —
  os coeficientes correspondem exatamente à ordem das variáveis da fórmula.
\end{nota}

\subsection*{J-test rejeita — instrumento inválido}

Se $z_2$ for contaminado com $\varepsilon$, o teste J detecta a violação:

\begin{lstlisting}[caption={Instrumento inválido — J-test rejeita}]
// z2b correlacionado com eps: instrumento inválido
generate df z2b = z2 + 0.5*eps
let m_bad = gmm(y ~ x, ~ z1 + z2b, df)
display m_bad
\end{lstlisting}

\begin{lstlisting}[language={}, basicstyle=\ttfamily\scriptsize,
  frame=none, numbers=none, backgroundcolor=\color{codbg},
  xleftmargin=0pt, xrightmargin=0pt, breaklines=false]
Dep. Variable: y  GMM Two-Step  N=1000  J=152.9417  Prob(J)=0.0000  DF=1
x0   1.3738***  (0.0089)    x1   2.2125***  (0.0124)
\end{lstlisting}

$J = 152{,}94$ ($p = 0{,}000$) rejeita $H_0$ com alta significância — e os
coeficientes são viesados ($\hat\beta_1 = 2{,}21$). O teste J funcionou como
alarme: antes de confiar nas estimativas GMM, verifique se $J$ não rejeita.

% ─────────────────────────────────────────────────────────────────────────────
\section{Sintaxe}

\begin{lstlisting}
gmm(y ~ x1 + x2, ~ z1 + z2 + z3, df)
gmm(y ~ x1 + x2, ~ z1 + z2 + z3, df, cov=robust)
\end{lstlisting}

O primeiro lado direito (\lstinline{~ x1 + x2}) lista os regressores;
o segundo (\lstinline{~ z1 + z2 + z3}) lista os instrumentos. Variáveis
exógenas em $x$ devem aparecer também como instrumentos de si mesmas.

\begin{lstlisting}[caption={Fluxo típico com variável exógena inclusa}]
load "dados.csv" as df

// edu é exógena (instrumento de si mesma)
// exp é endógena; z1 e z2 são instrumentos para exp
let m = gmm(sal ~ edu + exp, ~ edu + z1 + z2, df)
display m
\end{lstlisting}

\section{Capturando o resultado}

\begin{lstlisting}
let m = gmm(y ~ x, ~ z1 + z2, df)
display m
let yhat = predict(m, "xb")
\end{lstlisting}
